% SEGMENT OLP-0047-S01
% Part: sets-functions-relations
% Chapter: arithmetization
% Section: checking-details
%
\documentclass[../../../include/open-logic-section]{subfiles}


\begin{document}
	
\olfileid{sfr}{arith}{check}
\olsection{வரிசைப்படுத்தப்பட்ட வளையங்களும் புலங்களும்}

இந்த அத்தியாயம் முழுவதிலும், சில வரையறைகள் ``செயல்பட வேண்டியவாறே''
செயல்படுகின்றன என்று கூறினோம்.
இந்தத் தொழில்நுட்ப இணைப்பில், அதனால் நாம் குறிப்பது என்ன என்பதை
விரித்துரைத்து, வரையறைகள் ``சரியாகச்'' செயல்படுகின்றன என
(எவ்வாறு நிறுவுவது என்பதன் வரைவைக்) காட்டுவோம்.

\olref[int]{sec} இல், $\Int$ மீது கூட்டலையும் பெருக்கலையும்
வரையறுத்தோம். அவை வரையறுக்கப்பட்ட விதத்தில்,
$\Int$ க்கு நாம் ``விரும்பும்'' அமைப்பைத் தருகின்றன எனக் காட்ட வேண்டும்.
குறிப்பாக, கேள்விக்குரிய அமைப்பு ஒரு பரிமாற்று வளையத்தின் அமைப்பாகும்.

% SEGMENT OLP-0047-S02
\begin{defn}
	ஒரு \emph{பரிமாற்று வளையம்} என்பது, குறிப்பிட்ட $0$, $1$ என்ற
    உறுப்புகளையும் $+$, $\times$ என்ற செயல்களையும் கொண்ட ஒரு கணம் $S$;
    அது பின்வரும் எட்டு வாய்பாடுகளை நிறைவேற்றுகிறது:
	\begin{align*}
		\emph{சேர்ப்புப் பண்பு}&&a + (b+ c) & = (a + b) + c \\
		&& (a \times b) \times c & = a \times (b\times c)\\
		\emph{பரிமாற்றுப் பண்பு}&&a + b &= b+ a  \\
		&&  a \times b&= b\times a\\
		\emph{சமனிகள்}&&a + 0 &= a \\
		&& a \times 1 &= a\\
		\emph{கூட்டல் நேர்மாறு}&&(\exists b\in S)0&=a + b\\
		\emph{பங்கீட்டுப் பண்பு}&&a \times (b+ c ) &= (a \times b) + (a \times c)
	\end{align*}
	மறைமுகமாக, இவை அனைத்தும் $S$ க்கு வரம்பிடப்பட்ட அனைத்துப்
    பொதுநிலையளவிகளால் பிணைக்கப்பட்டுள்ளன.
    மேலும், இங்குள்ள $0$, $1$ என்ற உறுப்புகள்,
    அதே பெயருள்ள இயல் எண்களாக இருக்க வேண்டியதில்லை.
\end{defn}

% SEGMENT OLP-0047-S03
எனவே முழுக்கள் ஒரு பரிமாற்று வளையத்தை உருவாக்குகின்றனவா எனச்
சரிபார்க்க, இந்த எட்டு நிபந்தனைகளும் நிறைவேறுகின்றனவா என்று மட்டும்
சரிபார்க்க வேண்டும்.
எந்த நிபந்தனையையும் நிறுவுவது {கடினம்} அல்ல; ஆனால் இது சற்று உழைப்பானது.
எடுத்துக்காட்டாக, கூட்டலுக்கான \emph{சேர்ப்புப் பண்பை}
எவ்வாறு நிறுவுவது என்பது இதோ:

\begin{proof}
$i, j, k \in \Int$ என்பதை நிலைப்படுத்துக.
அப்போது $i = \equivrep{a_1, b_1}{}$, $j =
\equivrep{a_2,b_2}{}$, $k = \equivrep{a_3, b_3}{}$ என அமையும்
$a_1, b_1, a_2, b_2, a_3, b_3 \in \Nat$ உள்ளன.
(எளிதாகப் படிப்பதற்காக ``$\equivrep{x, y}{\Intequiv}$'' என்பதற்குப்
பதிலாக ``$\equivrep{x, y}{}$'' என எழுதுகிறோம்;
இந்தப் பிரிவு முழுவதிலும் இவ்வாறு செய்வோம்.) இப்போது:
\begin{align*}
	i + (j + k) &= \equivrep{a_1, b_1}{}+(\equivrep{a_2, b_2}{} + \equivrep{a_3, b_3}{}) \\
	&= \equivrep{a_1,  b_1}{} + \equivrep{a_2+a_3, b_2+b_3}{}\\
	&= \equivrep{a_1 + (a_2 + a_3), b_1 + (b_2 + b_3)}{}\\
	&= \equivrep{(a_1 + a_2) + a_3, (b_1 + b_2) + b_3}{}\\
	&= \equivrep{a_1 + a_2, b_1 + b_2}{} + \equivrep{a_3, b_3}{}\\
	&= (\equivrep{a_1, b_1}{} + \equivrep{a_2, b_2}{}) + \equivrep{a_3, b_3}{}\\
	&= (i+j) + k
\end{align*}
இங்கு $\Nat$ மீதான கூட்டலின் செயல்பாட்டைத் தடையின்றிப் பயன்படுத்தினோம்.
\end{proof}

% SEGMENT OLP-0047-S04
அதேபோல, \emph{கூட்டல் நேர்மாறை} எவ்வாறு நிறுவுவது என்பது இதோ:

\begin{proof}
$i \in \Int$ என்பதை நிலைப்படுத்துக;
ஏதோ $a,b \in \Nat$ க்கு $i = \equivrep{a,b}{}$.
$j = \equivrep{b,a}{} \in \Int$ என்க.
இயல் எண்களின் செயல்பாட்டைப் பயன்படுத்தினால்,
$(a+b) + 0 = 0 + (a+b)$; எனவே வரையறையின்படி
$\tuple{a+b, b+a} \sim_\Int \tuple{0,0}$;
ஆகவே $\equivrep{a+b, b+a}{} = \equivrep{0, 0}{} = 0_\Int$.
இப்போது $i + j =
\equivrep{a,b}{}+\equivrep{b,a}{}=\equivrep{a+b, b+a}{}= \equivrep{0,
0}{} = 0_\Int$.
\end{proof}

% SEGMENT OLP-0047-S05
\emph{பங்கீட்டுப் பண்பிற்கான} ஒரு நிறுவல் இதோ:

\begin{proof}
மேலே செய்ததுபோல, $i = \equivrep{a_1, b_1}{}$, $j =
\equivrep{a_2,b_2}{}$, $k = \equivrep{a_3, b_3}{}$ என
நிலைப்படுத்துக. இப்போது:
\begin{align*}
	i \times (j + k)
	&= \equivrep{a_1, b_1}{} \times (\equivrep{a_2,b_2}{} + \equivrep{a_3, b_3}{})\\
	&= \equivrep{a_1, b_1}{} \times \equivrep{a_2 + a_3,b_2+b_3}{}\\
	&= \equivrep{a_1  (a_2 + a_3) + b_1  (b_2+b_3), a_1  (b_2 + b_3) + b_1 (a_2 + a_3)}{}\\
	&= \equivrep{a_1 a_2 + a_1a_3 + b_1 b_2+b_1b_3, a_1 b_2 + a_1b_3 + a_2b_1 + a_3b_1}{}\\		
%		&= \equivrep{a_1 a_2 + b_1b_2 + a_1a_3 + b_1 b_2+b_1b_3, a_1 b_2 + a_1b_3 + a_2b_1 + a_3b_1}{}\\		
	&= \equivrep{a_1a_2 + b_1b_2, a_1b_2 + a_2b_1}{} + \equivrep{a_1a_3 + b_1b_3, a_1b_3 + a_3b_1}{}\\
	&= (\equivrep{a_1, b_1}{} \times \equivrep{a_2,b_2}{}) + (\equivrep{a_1, b_1}{} \times  \equivrep{a_3, b_3}{})\\
	&= (i \times j) + (i \times k)
\end{align*}
\end{proof}

% SEGMENT OLP-0047-S06
மீதமுள்ள ஐந்து நிபந்தனைகளை நிறுவுவதை ஒரு பயிற்சியாக விடுகிறோம்.
அவற்றைச் செய்தபின், $\Int$ ஒரு பரிமாற்று வளையத்தை அமைக்கிறது,
அதாவது கூட்டலும் பெருக்கலும் (வரையறுக்கப்பட்டவாறு)
செயல்பட வேண்டியவாறே செயல்படுகின்றன, எனக் காட்டியிருப்போம்.

\begin{prob}
$\Int$ ஒரு பரிமாற்று வளையம் என நிறுவுக.
\end{prob}

% SEGMENT OLP-0047-S07
ஆனால் நமது பணி முடிந்துவிடவில்லை.
$\Int$ மீது கூட்டலையும் பெருக்கலையும் வரையறுத்ததோடு,
$\leq$ என்ற ஒரு வரிசைத் தொடர்பையும் வரையறுத்தோம்;
இதுவும் செயல்பட வேண்டியவாறு செயல்படுகிறதா எனச் சரிபார்க்க வேண்டும்.
மேலும் விரிவாக, $\Int$ ஒரு \emph{வரிசைப்படுத்தப்பட்ட} வளையத்தை
அமைக்கிறது எனக் காட்ட வேண்டும்.
\footnote{\olref[sfr][rel][ord]{def:linearorder} இலிருந்து நினைவுகொள்க:
	ஒரு முழு வரிசை என்பது தற்சுட்டு, கடப்பு, எதிர்சமச்சீர்,
    ஒப்பிடத்தக்க ஆகிய பண்புகளைக் கொண்ட தொடர்பாகும்.
	வரிசைத் தொடர்புகளின் சூழலில், ஒப்பிடத்தக்க பண்பு சில நேரங்களில்
	\emph{முக்கூற்றுப் பண்பு} எனப்படும்; ஏனெனில் எந்த $a$, $b$ க்கும்
    $a \leq b \lor a = b \lor a \geq b$.}

% SEGMENT OLP-0047-S08
\begin{defn}
ஒரு \emph{வரிசைப்படுத்தப்பட்ட வளையம்} என்பது,
$\leq$ என்ற முழு வரிசைத் தொடர்பையும் கொண்ட ஒரு பரிமாற்று வளையம்;
அது பின்வருவனவற்றை நிறைவேற்றுகிறது:
\begin{align*}
	a \leq b &\lif a + c \leq b + c\\
	(a \leq b \land 0 \leq c) &\lif a \times c \leq b \times c
\end{align*}
\end{defn}

% SEGMENT OLP-0047-S09
\begin{prob}
$\Int$ ஒரு வரிசைப்படுத்தப்பட்ட வளையம் என நிறுவுக.
\end{prob}

% SEGMENT OLP-0047-S10
முன்புபோல, கட்டமைக்கப்பட்ட $\Int$ ஒரு வரிசைப்படுத்தப்பட்ட வளையம்
எனக் காட்டுவது உழைப்பானது, ஆனால் வழக்கமானது.
அதனை உங்களுக்கே விட்டுவிடுகிறோம்.

இது முழுக்களைக் கவனித்துக்கொள்கிறது.
ஆனால் இப்போது விகிதமுறு எண்களைப் பற்றி இதே போன்றவற்றைக் காட்ட வேண்டும்.
குறிப்பாக, $+$, $\times$, $\leq$ ஆகியவற்றுக்கு நாம் கொடுத்த
வரையறைகளின் கீழ், விகிதமுறு எண்கள் ஒரு வரிசைப்படுத்தப்பட்ட
\emph{புலத்தை} அமைக்கின்றன எனக் காட்ட வேண்டும்:
\begin{defn}\ollabel{orderedfield}
ஒரு \emph{வரிசைப்படுத்தப்பட்ட புலம்} என்பது,
பின்வருவதையும் நிறைவேற்றும் வரிசைப்படுத்தப்பட்ட வளையமாகும்:
\begin{align*}
	\emph{பெருக்கல் நேர்மாறு}& & (\forall a \in S \setminus \{0\})(\exists b \in S) a\times b& = 1
\end{align*}
\end{defn}

% SEGMENT OLP-0047-S11
$\Int$ ஒரு வரிசைப்படுத்தப்பட்ட வளையம் என்று காட்டியவுடன்,
$\Rat$ ஒரு வரிசைப்படுத்தப்பட்ட புலம் என்று காட்டுவது எளிதானது,
ஆனால் உழைப்பானது.

\begin{prob}
$\Rat$ ஒரு வரிசைப்படுத்தப்பட்ட புலம் என நிறுவுக.
\end{prob}

% SEGMENT OLP-0047-S12
முழுக்களையும் விகிதமுறு எண்களையும் கவனித்தபின்,
மெய்யெண்களை மட்டும் கவனிக்க வேண்டியுள்ளது.
குறிப்பாக, $\Real$ ஒரு \emph{முழுமையான} வரிசைப்படுத்தப்பட்ட புலம்,
அதாவது முழுமைப் பண்பைக் கொண்ட வரிசைப்படுத்தப்பட்ட புலம்,
என்று காட்ட வேண்டும்.
\olref[cuts]{realcompleteness} ஆனது $\Real$ முழுமைப் பண்பைக்
கொண்டுள்ளது என நிறுவியது.
ஆயினும், $\Real$ ஒரு வரிசைப்படுத்தப்பட்ட புலம் என்பதைச் சரிபார்க்கும்
(சலிப்பூட்டும்) பணியை இன்னும் முழுவதும் செய்ய வேண்டும்.

\emph{அந்த} உழைப்பான பயிற்சியை நோக்கி விரைவதற்கு முன்,
மேலும் சில ``உடனடியான'' உண்மைகளைச் சரிபார்க்க வேண்டும்.
எடுத்துக்காட்டாக, எந்த வெட்டுகள் $\alpha$, $\beta$ க்கும்,
வரையறுக்கப்பட்ட $\alpha + \beta$ உண்மையில் ஒரு \emph{வெட்டு}
என்பதற்கான உறுதி தேவை. அந்த உண்மைக்கான நிறுவல் இதோ:

% SEGMENT OLP-0047-S13
\begin{proof}	
$\alpha$, $\beta$ இரண்டும் வெட்டுகள் என்பதால்,
$\alpha + \beta = \Setabs{p
+ q}{p \in \alpha \land q \in \beta}$ என்பது $\Rat$ இன்
வெற்றற்ற தகு உட்கணம்.
இப்போது ஏதோ $p \in \alpha$, $q \in \beta$ க்கு
$x < p + q$ எனக் கொள்க.
அப்போது $x - p < q$; எனவே $x - p \in \beta$;
மேலும் $x = p + (x - p) \in \alpha + \beta$.
ஆகவே $\alpha + \beta$ என்பது $\Rat$ இன் ஒரு தொடக்கத் துண்டு.
இறுதியாக, எந்த $p + q \in \alpha + \beta$ க்கும்,
$\alpha$, $\beta$ இரண்டும் வெட்டுகள் என்பதால்,
$p < p_1$, $q < q_1$ என அமையும் $p_1 \in \alpha$, $q_1 \in \beta$
உள்ளன; எனவே $p + q < p_1 + q_1 \in \alpha +
\beta$; ஆகவே $\alpha + \beta$ க்கு மீப்பெரு உறுப்பு இல்லை.
\end{proof}

% SEGMENT OLP-0047-S14
இதே போன்ற முயற்சிகள், $\alpha - \beta$, $\alpha \times \beta$,
$\alpha \div \beta$ ஆகியவை வெட்டுகள் எனச் சரிபார்க்க உதவும்
(இறுதி நிலையில், $\beta$ பூச்சிய வெட்டாக இருக்கும் நிலையைத் தவிர்க்கவும்).
ஆனால் இதையும் உங்களுக்கே விட்டுவிடுகிறோம்.

\begin{prob}
$\Real$ ஒரு வரிசைப்படுத்தப்பட்ட புலம் என நிறுவுக.
\end{prob}

% SEGMENT OLP-0047-S15
சரிசெய்ய வேண்டிய ஒரு சிறிய எஞ்சிய பகுதி உள்ளது.
\olref[cuts]{sec} இல், $\sqrt{2} =
\Setabs{p \in \Rat}{p < 0 \text{ அல்லது }p^2 < 2}$ எனக் கொள்ளலாம்
என்று முன்மொழிகிறோம்.
ஆனால் இந்தக் கணம் ஒரு \emph{வெட்டு} என்று காட்ட வேண்டும்.
அந்த உண்மைக்கான நிறுவல் இதோ:

% SEGMENT OLP-0047-S16
\begin{proof}
இது விகிதமுறு எண்களின் வெற்றற்ற தகு தொடக்கத் துண்டு என்பது தெளிவு;
எனவே இதற்கு மீப்பெரு உறுப்பு இல்லை எனக் காட்டினால் போதும்.
குறிப்பாக, $p$ என்பது $p^2 < 2$ என அமையும் நேர்ம விகிதமுறு எண்ணாகவும்,
$q = \frac{2p+2}{p+2}$ ஆகவும் இருந்தால்,
$p < q$, $q^2 < 2$ இரண்டையும் காட்டினால் போதும்.
$p < q$ என்பதைக் காண, பின்வருவதை மட்டும் கவனிக்கவும்:
\begin{align*}
	p^2 &< 2\\
	p^2 + 2p &< 2 + 2p\\
	p(p + 2) &< 2 + 2p\\
	p &< \tfrac{2+2p}{p+2} = q
\end{align*}
$q^2 < 2$ என்பதைக் காண, பின்வருவதை மட்டும் கவனிக்கவும்:
\begin{align*}
	p^2 &< 2\\
	2p^2 + 4p + 2 &< p^2 + 4p+ 4\\
	4p^2 + 8p + 4 &< 2(p^2 + 4p + 4)\\
	(2p+2)^2 & <2(p+2)^2\\
	\tfrac{(2p+2)^2}{(p+2)^2} &< 2\\
	q^2 &< 2
\end{align*}
\end{proof}
\end{document}
